%%%%%
%
%% Dans ce fichier, les �ventuels accents sont cod�s
%% en ISO-8859-15 (aka latin9)
%% La ligne suivante donne trois exemples 
%% e accent aigu : �, e accent grave : �, e-dans-l'o : �
\section{Filtre num�rique �quivalent � un circuit RC sans d�phasage du signal de sortie sur des plages temporelles}

\itemstyle{Op\'erateur~:} \operfmt{Filtre\_RC\_sdph\_paliers}
\index{Filtre\_RC\_sdph\_paliers}

%%%%%OBJET%%%%
\itemstyle{Objet~:}
Cet op�rateur applique � un param\`etre un filtrage num�rique �quivalent
� un circuit �lectrique RC. Il est d�fini par sa fr�quence de coupure � -3dB $f_c$.

L'utilisateur sp�cifie si l'op�rateur doit appliquer un filtre passe-bas ou 
passe-haut avec l'argument $Type$.

Le filtrage est effectu� sur les plages temporelles d�finies en entr�e. En dehors de ces plages temporelles on trouvera des valeurs manquantes.


%%%%%ARGUMENTS%%%%
\bigskip
%% Dans ce fichier, les �ventuels accents sont cod�s
%% en ISO-8859-15 (aka latin9)
%% La ligne suivante donne trois exemples 
%% e accent aigu : �, e accent grave : �, e-dans-l'o : �
%%%%%ENTREES%%%%
\begin{tabarg}
\multicolumn{3}{l}{\titretab {Entr\'ees}}\\
\hline
$X$ & Vecteur, {\tt float} & Param�tre d'entr�e \\
$f_c$ & Coeff., {\tt float} & Fr�quence de coupure du filtre ($Hz$)\\
$Type$ & Coeff., {\tt Int\_32} & Type de filtre (0 : passe-bas, 1: passe-haut)\\
$time_out$ & Vecteur, {\tt float} & R�f�rence de temps qui d�finit les plages de calcul \\
\end{tabarg}

\medskip
%%%%%SORTIES%%%%
\begin{tabarg}
\multicolumn{3}{l}{\titretab {Sorties}}\\
\hline
$X_out$ & Vecteur, {\tt float} & Param\`etre filtr\'e (m\^eme unit\'e que $X$) \\
\end{tabarg}


%%%%%FORMULE%%%%
\itemstyle{Formule}

%% Formule de l'operateur : Filtre_RC_sdph_paliers
%% Dans ce fichier, les �ventuels accents sont cod�s
%% en ISO-8859-15 (aka latin9)
%% La ligne suivante donne trois exemples 
%% e accent aigu : �, e accent grave : �, e-dans-l'o : �
\[
X_f[i] = e^{-\frac{2\pi.f_c}{f_e}}.X_f[i-1] + 
         \left( 1 - e^{-\frac{2\pi.f_c}{f_e}} \right). X[i]
\]
si $Type = 1$ alors $X_f = X - X_f$



%%%%%AUTEUR%%%%
\itemstyle{Auteur~:}
CNRM/GMEI/TRAMM

%%%%%REFERENCE%%%%
\itemstyle{R\'ef\'erences~:}
